\documentclass{article}
\usepackage[margin=1in]{geometry}
\usepackage{amsmath,amssymb,amsthm}
\usepackage{booktabs,longtable}
\usepackage{hyperref}

\newtheorem{definition}{Definition}
\newtheorem{theorem}[definition]{Theorem}
\newtheorem{lemma}[definition]{Lemma}
\newtheorem{proposition}[definition]{Proposition}
\newtheorem{corollary}[definition]{Corollary}

\newcommand{\source}[1]{\par\noindent\textbf{Source.} #1}
\newcommand{\lean}[1]{\par\noindent\textbf{Lean target.} \texttt{\detokenize{#1}}}
\newcommand{\status}[1]{\par\noindent\textbf{Status.} \texttt{\detokenize{#1}}}
\newcommand{\depends}[1]{\par\noindent\textbf{Depends on.} #1}
\newcommand{\cM}{\mathcal M}
\newcommand{\cZ}{\mathcal Z}
\newcommand{\cH}{\mathcal H}
\newcommand{\cB}{\mathcal B}

\begin{document}

\section*{Blueprint 011: Project-Local Von Neumann Type Theory}

\subsection*{Scope}

This note promotes the fourth remaining blueprint block from
\path{agents/2026-06-22-remaining-blueprint-suite-goal.md}.  It supports
\[
\texttt{VN-002},\quad
\texttt{VECTOR-ANALYTIC-001},\quad
\texttt{PROJ-001},\quad
\texttt{CENTRAL-SUPPORT-001},\quad
\texttt{TRACE-001},\quad
\texttt{TYPE-DEF-001},\quad
\texttt{TYPE-ALT-001},\quad
\texttt{MATRIX-001},
\]
and the supporting represented-region criteria
\[
\texttt{REGION-VECTOR-001},\quad
\texttt{REGION-VECTOR-NORMAL-001},\quad
\texttt{REGION-VECTOR-TRACE-001},\quad
\texttt{REGION-VECTOR-FAITH-001},\quad
\texttt{TRACE-TYPE-CRITERION-001},\quad
\texttt{REGION-VECTOR-TYPE-001}.
\]
It gives the project-local vocabulary for centers, factors, projection
comparison, traces, factor types, and the vector-state criteria used later for
represented region algebras.  It does not prove the final stabilizer
classification and does not create Aristotle packets; the Lean-facing names
below are blueprint targets, not integrated Lean declarations.
The represented-region IDs above are mirrored in \path{MANIFEST.toml} and
remain pending adversarial review and formalization.

\subsection*{Takesaki Statement Transcripts}

The local Takesaki I PDF is scanned; the cited Chapter V pages were rendered
and OCR-checked in \path{/tmp/formalization-ocr}.

\begin{definition}[VonNeumannAlgebraFactor, Takesaki I Definition II.3.2, PDF p.~92]
A von Neumann algebra on \(\mathcal H\) is a \(*\)-subalgebra
\(\mathcal M\) of \(\mathcal L(\mathcal H)\) such that
\[
  \mathcal M=\mathcal M''.
\]
A factor is a von Neumann algebra \(\mathcal M\) with trivial center,
\(\mathcal M\cap\mathcal M'=\mathbb C1\).
\end{definition}

\begin{definition}[ProjectionEquivalence, Takesaki I Definition V.1.2, PDF p.~310]
Two projections \(e\) and \(f\) in a von Neumann algebra \(\mathcal M\) are
said to be equivalent if there exists an element \(u\in\mathcal M\) with
\[
  u^*u=e,\qquad uu^*=f.
\]
This is written \(e\sim f\).  The projections \(e\) and \(f\) are called,
respectively, the initial projection and the final projection of \(u\).  If a
projection \(e\in\mathcal M\) is equivalent to a projection \(f_0\in\mathcal M\)
with \(f_0\leq f\in\mathcal M\), then one writes \(e\precsim f\).
\end{definition}

\begin{theorem}[ProjectionComparability, Takesaki I Theorem V.1.8, PDF p.~313]
For any pair \(e,f\) of projections in a von Neumann algebra \(\mathcal M\),
there exists a central projection \(z\) such that
\[
  ze\precsim zf,\qquad (1-z)e\succeq (1-z)f.
\]
In particular, if \(\mathcal M\) is a factor, then one and only one of the
following relations holds:
\[
  e\prec f,\qquad e\sim f,\qquad e\succ f.
\]
\end{theorem}

\begin{proof}
Choose a maximal family of orthogonal subprojections \(e_i\leq e\) and
\(f_i\leq f\) with \(e_i\sim f_i\), and let \(e_0=\sum_i e_i\) and
\(f_0=\sum_i f_i\).  Maximality and the central-orthogonality lemma imply
that \(e-e_0\) and \(f-f_0\) are separated by a central projection \(z\).
Multiplying the equivalence \(e_0\sim f_0\) by \(z\) and \(1-z\) gives the
two comparison inequalities.  In a factor, the only central projections are
\(0\) and \(1\), giving the trichotomy.
\end{proof}

\begin{definition}[FiniteInfiniteProjection, Takesaki I Definition V.1.15, PDF p.~316]
A projection \(e\) in a von Neumann algebra \(\mathcal M\) is said to be
finite if \(e\sim f<e\) implies \(e=f\).  Otherwise it is said to be
infinite.  A projection \(e\) is said to be purely infinite if there is no
nonzero finite projection \(f<e\) in \(\mathcal M\).  If \(ze\) is infinite
for every central projection \(z\in\mathcal M\) with \(ze\neq0\), then \(e\)
is called properly infinite.  If \(e\mathcal M e\) is abelian, then \(e\) is
said to be abelian.
\end{definition}

\begin{definition}[FiniteInfiniteVonNeumannAlgebra, Takesaki I Definition V.1.16, PDF p.~316]
A von Neumann algebra \(\mathcal M\) is said to be finite, infinite, properly
infinite, or purely infinite according to the property of the identity
projection \(1\).
\end{definition}

\begin{definition}[VonNeumannAlgebraTypes, Takesaki I Definition V.1.17, PDF p.~316]
A von Neumann algebra \(\mathcal M\) is said to be of type I if every nonzero
central projection in \(\mathcal M\) majorizes a nonzero abelian projection
in \(\mathcal M\).  If there is no nonzero finite projection in
\(\mathcal M\), that is, if \(\mathcal M\) is purely infinite, then it is said
to be of type III.  If \(\mathcal M\) has no nonzero abelian projection and if
every nonzero central projection in \(\mathcal M\) majorizes a nonzero finite
projection of \(\mathcal M\), then it is said to be of type II.  If
\(\mathcal M\) is finite and of type II, then it is said to be of type
\(II_1\).  If \(\mathcal M\) is of type II and has no nonzero central finite
projection, then \(\mathcal M\) is said to be of type \(II_\infty\).
\end{definition}

\begin{theorem}[CentralTypeDecomposition, Takesaki I Theorem V.1.19, PDF pp.~316--317]
Every von Neumann algebra \(\mathcal M\) is uniquely decomposable into the
direct sum of those of type I, type \(II_1\), type \(II_\infty\), and type
III.  Moreover, every projection \(e\) in \(\mathcal M\) is uniquely written
as the sum of centrally orthogonal projections \(e_1\) and \(e_2\) in
\(\mathcal M\) such that \(e_1\) is finite and \(e_2\) is properly infinite.
\end{theorem}

\begin{proof}
Take a maximal centrally orthogonal family of abelian projections; its central
support gives the type I summand.  In the complement, take a maximal centrally
orthogonal family of finite projections; its central support gives the type II
summand and the remaining central part is type III.  Splitting the type II
part by maximal central finite projections separates \(II_1\) and
\(II_\infty\).  The same central-support argument proves uniqueness.  Applying
the construction inside the corner \(e\mathcal M e\) gives the projection
decomposition.
\end{proof}

\begin{theorem}[FactorTypeAlternatives, Takesaki I Corollary V.1.20, PDF p.~317]
A factor is either of type I, type \(II_1\), type \(II_\infty\), or type III.
\end{theorem}

\begin{proof}
Apply Theorem V.1.19.  Since the center of a factor has no nontrivial
projections, exactly one central type summand can be nonzero.
\end{proof}

\begin{definition}[SemifiniteVonNeumannAlgebra, Takesaki I Definition V.1.21, PDF p.~317]
In the decomposition of Theorem V.1.19, if the type III central summand is
zero, then \(\mathcal M\) is said to be semifinite.
\end{definition}

\begin{proposition}[MatrixAlgebraOverAbelianTypeI, Takesaki I Proposition V.1.23, PDF p.~318]
\leavevmode
\begin{enumerate}
\item If \(\mathcal A\) is an abelian von Neumann algebra, then
\(\mathcal A\bar\otimes\mathcal L(\mathcal H)\) is of type I for any Hilbert
space \(\mathcal H\).
\item If \(\mathcal A\) is an abelian von Neumann algebra, then
\(M_n(\mathcal A)\), the \(n\times n\) matrix algebra with entries from
\(\mathcal A\), is a finite von Neumann algebra of type I, and the
\(\mathcal A\)-valued map
\[
  \Phi(x)=\frac1n\sum_i x_{i,i},
  \qquad x=[x_{i,j}]\in M_n(\mathcal A),
\]
satisfies
\[
  \Phi(a)=a,\qquad a\in\mathcal A,
\]
\[
  \Phi(axb)=a\Phi(x)b,\qquad a,b\in\mathcal A,
\]
and
\[
  \Phi(x^*x)=\Phi(xx^*)>0,\qquad
  \Phi(x^*x)=0\Rightarrow x=0.
\]
\end{enumerate}
\end{proposition}

\begin{proof}
For the tensor product statement, take a minimal projection \(p\) in
\(\mathcal L(\mathcal H)\); then \(1\otimes p\) is abelian and has central
support \(1\).  For the matrix statement, the displayed properties of
\(\Phi\) follow from the normalized diagonal sum and direct matrix
multiplication, and faithfulness follows from summing the positive diagonal
entries of \(x^*x\).
\end{proof}

\begin{theorem}[TypeIFactorIsBoundedOperators, Takesaki I Corollary V.1.28, PDF p.~320]
A factor of type I is isomorphic to \(\mathcal L(\mathcal H)\) for some
Hilbert space \(\mathcal H\).  Therefore, the dimension of
\(\mathcal H\) determines completely the algebraic type of the factor.
\end{theorem}

\begin{proof}
In a factor of type I, the identity is homogeneous of some cardinal dimension.
The preceding homogeneous type I structure theorem identifies the algebra with
operators on a Hilbert space of that dimension.
\end{proof}

\begin{definition}[TypeICardinalHomogeneous, Takesaki I Definition V.1.33, PDF p.~322]
A von Neumann algebra \(\mathcal M\) of type I is said to be of type
\(I_\alpha\) if \(1\) is \(\alpha\)-homogeneous.  By Corollary V.1.28, a
factor of type \(I_n\), \(n<+\infty\), is nothing but the \(n\times n\) matrix
algebra over \(\mathbb C\).
\end{definition}

\begin{definition}[TraceOnVonNeumannAlgebra, Takesaki I Definition V.2.1, PDF p.~329]
A trace on a von Neumann algebra \(\mathcal M\) is a function \(\tau\) on the
positive cone \(\mathcal M_+\), with values in the extended positive reals
\([0,+\infty]\), satisfying
\[
  \tau(x+y)=\tau(x)+\tau(y),\qquad x,y\in\mathcal M_+,
\]
\[
  \tau(\lambda x)=\lambda\tau(x),\qquad \lambda\geq0,\ x\in\mathcal M_+,
\]
and
\[
  \tau(x^*x)=\tau(xx^*),\qquad x\in\mathcal M,
\]
with the usual convention \(0(+\infty)=0\).  A trace \(\tau\) is said to be
faithful if \(\tau(x)>0\) for any nonzero \(x>0\), semifinite if every nonzero
\(x\in\mathcal M_+\) majorizes some nonzero \(y>0\) with
\(\tau(y)<+\infty\), finite if \(\tau(1)<+\infty\), and normal if
\[
  \tau(\sup x_i)=\sup\tau(x_i)
\]
for every bounded increasing net \(\{x_i\}\) in \(\mathcal M_+\).
\end{definition}

\begin{theorem}[FiniteIffEnoughFiniteNormalTraces, Takesaki I Theorem V.2.4, PDF p.~331]
For a von Neumann algebra \(\mathcal M\), the following conditions are
equivalent:
\begin{enumerate}
\item \(\mathcal M\) is finite;
\item \(\mathcal M\) admits sufficiently many finite normal traces.
\end{enumerate}
\end{theorem}

\begin{proof}
If sufficiently many finite traces exist, an isometry \(u\) with \(u^*u=1\)
satisfies \(\tau(uu^*)=\tau(1)\) for every such trace, so \(1-uu^*=0\).
Conversely, for a finite algebra, average normal positive functionals over
unitary conjugacy using the fixed-point theorem to obtain finite normal
traces.  Their supports cover the identity because their central supports
match the supports of the original normal functionals on the center.
\end{proof}

\begin{theorem}[SemifiniteIffFaithfulSemifiniteNormalTrace, Takesaki I Theorem V.2.15, PDF p.~337]
For a von Neumann algebra \(\mathcal M\), the following conditions are
equivalent:
\begin{enumerate}
\item \(\mathcal M\) is semifinite;
\item \(\mathcal M\) admits a faithful semifinite normal trace.
\end{enumerate}
\end{theorem}

\begin{proof}
A faithful semifinite normal trace makes every nonzero projection majorize a
nonzero finite projection, by cutting a finite-trace positive element by a
nonzero spectral projection.  Conversely, take a maximal family of semifinite
normal traces with mutually orthogonal supports.  If the total support were
not \(1\), the remaining semifinite central summand would contain a nonzero
part of the form \(\mathcal N\bar\otimes\mathcal L(\mathcal H)\) with
\(\mathcal N\) finite; a finite normal trace on \(\mathcal N\) and the
standard trace construction would contradict maximality.
\end{proof}

\subsection*{Centers and Factors}

\begin{definition}[CenterAndFactor, ID: VN-002, Takesaki I Definition II.3.2, PDF p.~92]
Let \(\cM\subseteq\cB(\cH)\) be a concrete von Neumann algebra.  Its center is
\[
  \cZ(\cM)=\cM\cap\cM'.
\]
The algebra \(\cM\) is a factor if \(\cZ(\cM)=\mathbb C\,1\).
\source{Takesaki I, Definition II.3.2, PDF p.~92.}
\depends{\texttt{VN-001}.}
\lean{Formalization.OperatorAlgebra.Factor.center; Formalization.OperatorAlgebra.Factor.IsFactor}
\status{blueprint_note_complete}
\end{definition}

\subsection*{Projection Comparison}

\begin{definition}[ConcreteProjectionVocabulary, ID: PROJ-001, Takesaki I Definitions V.1.2 and V.1.15, PDF pp.~310, 316]
Let \(\cM\) be a concrete \(\mathrm{VonNeumannAlgebra}\,\cH\), viewed as a
unital \(*\)-subalgebra of \(\cB(\cH)\).
\begin{itemize}
\item A projection in \(\cM\) is an operator \(p\in\cM\) satisfying
\[
  p^*=p,\qquad p^2=p.
\]
\item For projections \(p,q\in\cM\), the projection order \(p\le q\) means
\[
  pq=p,
\]
equivalently \(qp=p\), equivalently \(\operatorname{Ran}(p)\subseteq
\operatorname{Ran}(q)\).  A subprojection of \(q\) is a projection \(p\in\cM\)
with \(p\le q\).  A proper subprojection is a subprojection \(p\le q\) with
\(p\neq q\).
\item A central projection is a projection \(z\in\cM\) lying in
\(\cZ(\cM)\); equivalently \(zx=xz\) for every \(x\in\cM\).
\item A partial isometry in \(\cM\) is an operator \(v\in\cM\) such that
\[
  v^*v\quad\hbox{and}\quad vv^*
\]
are projections in \(\cM\).  These are the initial and final projections of
\(v\).
\item For a projection \(e\in\cM\), the corner \(e\cM e\) is the set
\[
  e\cM e=\{exe:x\in\cM\}.
\]
It is treated as a unital von Neumann algebra with identity \(e\), represented
on the Hilbert subspace \(e\cH=\operatorname{Ran}(e)\).  A projection \(e\) is
abelian when this corner is commutative.
\end{itemize}
The well-formedness facts used by Lean are: \(p^*=p\) and \(p^2=p\) imply that
\(\operatorname{Ran}(p)\) is closed; if \(e\in\cM\), then \(e\cM e\) is closed
under multiplication and adjoint because \(\cM\) is; and strong-operator
closedness of \(e\cM e\) follows by compressing strong limits by the bounded
operator \(e\).
\source{Takesaki I, Definition V.1.2, PDF p.~310; Definition V.1.15,
PDF p.~316.}
\depends{\texttt{VN-002}.}
\lean{Formalization.OperatorAlgebra.ProjectionComparison}
\status{blueprint_note_complete}
\end{definition}

\begin{definition}[MurrayVonNeumannAndFinitePredicates, supporting ID: PROJ-001, Takesaki I Definitions V.1.2 and V.1.15--V.1.16, PDF pp.~310, 316]
For projections \(e,f\in\cM\), define Murray--von Neumann equivalence
\[
  e\sim_{\cM} f
\]
to mean that there is a partial isometry \(u\in\cM\) with
\[
  u^*u=e,\qquad uu^*=f.
\]
Define subequivalence \(e\precsim_{\cM} f\) to mean that there is a
projection \(f_0\in\cM\) such that \(f_0\le f\) and \(e\sim_{\cM} f_0\).  Write
\(e\prec_{\cM} f\) when \(e\precsim_{\cM} f\) but \(e\not\sim_{\cM} f\).

A projection \(e\in\cM\) is finite if every projection \(f\in\cM\) with
\[
  f\le e,\qquad f\sim_{\cM} e
\]
satisfies \(f=e\).  It is infinite if it is not finite.  It is properly
infinite if for every central projection \(z\in\cM\),
\[
  ze\neq0\quad\Longrightarrow\quad ze\hbox{ is infinite in }z\cM z.
\]
The algebra \(\cM\) is finite, infinite, or properly infinite when its
identity projection has the corresponding property.
\source{Takesaki I, Definition V.1.2, PDF p.~310; Definition V.1.15 and
Definition V.1.16, PDF p.~316.}
\depends{\texttt{VN-002}.}
\lean{Formalization.OperatorAlgebra.ProjectionComparison}
\status{blueprint_note_complete}
\end{definition}

\begin{definition}[CentralSupportProjectionLatticePackage, ID: CENTRAL-SUPPORT-001, Takesaki I Proposition V.1.1, Lemma V.1.7, Lemma V.1.18, Theorems V.1.8 and V.1.19, PDF pp.~310, 312--313, 316--317]
For Lean extraction, the central-support infrastructure used below is the
following source-backed package for a concrete von Neumann algebra \(\cM\).
It may be implemented either by proving these fields from the cited Takesaki
results or, temporarily, by importing them as explicitly named assumptions in
packets that need comparability or central decomposition.
\begin{itemize}
\item \(\operatorname{ProjSupInf}_{\cM}\): the projections of \(\cM\) form a
complete lattice.  For every family \((p_i)\), the projections
\(\bigvee_i p_i\) and \(\bigwedge_i p_i\) lie in \(\cM\).  If the \(p_i\)
are pairwise orthogonal, then \(\bigvee_i p_i\) is the strong-operator sum of
the \(p_i\).
\item \(\operatorname{CentralSupport}_{\cM}(p)\): for every projection
\(p\in\cM\), \(c_{\cM}(p)\) is the least central projection \(z\in\cM\) with
\(p\le z\).  For a central projection \(z\), central cutting satisfies
\[
  c_{z\cM z}(zp)=z\,c_{\cM}(p)
\]
inside the corner \(z\cM z\), and central cuts preserve projection order.
\item \(\operatorname{CentralCutMvN}_{\cM}\): if \(p\sim_{\cM}q\) is
implemented by a partial isometry \(v\in\cM\) and \(z\) is central, then
\[
  zp\sim_{z\cM z} zq
\]
is implemented by \(zv\).  Consequently central cuts preserve
subequivalence, finiteness, infiniteness, proper infiniteness, and abelian
corners.
\item \(\operatorname{CentralSupportSeparation}_{\cM}\): if projections
\(r,s\in\cM\) have no nonzero subprojections \(r_0\le r\) and \(s_0\le s\)
with \(r_0\sim_{\cM}s_0\), then \(c_{\cM}(r)c_{\cM}(s)=0\).  Equivalently,
there is a central projection separating \(r\) from \(s\).
\item \(\operatorname{OrthogonalPartialIsometrySum}_{\cM}\): for any family
of partial isometries \((v_i)\subseteq\cM\) with pairwise orthogonal initial
projections and pairwise orthogonal final projections, the strong sum
\(\sum_i v_i\) belongs to \(\cM\) and is a partial isometry whose initial and
final projections are \(\bigvee_i v_i^*v_i\) and \(\bigvee_i v_i v_i^*\).
\item \(\operatorname{CentrallyOrthogonalFiniteSup}_{\cM}\): suprema of
centrally orthogonal abelian projections are abelian, and suprema of
centrally orthogonal finite projections are finite.  In particular, maximal
centrally orthogonal families of abelian or finite projections have central
supports that define central type summands.
\item \(\operatorname{LargestCentralFiniteProjection}_{\cM}\): for every
central projection \(z\), the supremum of all central finite projections
below \(z\) exists and is the largest central finite subprojection of \(z\);
its complement \(z-z_f\) has no nonzero central finite subprojection.  This
is the construction later applied in \texttt{TYPE-ALT-001} to split the
central part that satisfies the type-II predicate.
\end{itemize}
This target is definition-only at the blueprint level: it names the exact
Lean-facing fields used by later proof routes and records their source
locators, instead of leaving central support, projection-lattice completeness,
orthogonal sums, or largest central finite projections implicit.
\source{Takesaki I, Proposition V.1.1, PDF p.~310, for the complete
projection lattice; Lemma V.1.7, PDF p.~312, for central support and central
orthogonality; Theorem V.1.8, PDF p.~313, for central cuts in comparability;
Lemma V.1.18 and Theorem V.1.19, PDF pp.~316--317, for centrally orthogonal
finite and abelian suprema and the largest central finite projection used in
the type decomposition.}
\depends{\texttt{VN-002}, \texttt{PROJ-001}.}
\lean{Formalization.OperatorAlgebra.CentralSupportProjectionLatticePackage}
\status{blueprint_note_complete}
\end{definition}

\begin{theorem}[ProjectionComparability, supporting ID: PROJ-001, Takesaki I Theorem V.1.8, PDF p.~313]
For projections \(e,f\) in a von Neumann algebra \(\cM\), there exists a
central projection \(z\in\cM\) such that
\[
  ze\precsim zf,\qquad
  (1-z)e\succeq (1-z)f.
\]
If \(\cM\) is a factor, exactly one of \(e\prec f\), \(e\sim f\), or
\(e\succ f\) holds.
\source{Takesaki I, Theorem V.1.8, PDF p.~313.}
\depends{\texttt{VN-002}, \texttt{PROJ-001}, \texttt{CENTRAL-SUPPORT-001}.}
\lean{Formalization.OperatorAlgebra.projection_comparability}
\status{blueprint_note_complete}
\end{theorem}

\begin{proof}
The formal proof route is the standard Murray--von Neumann comparison
argument.  Consider the partially ordered set of families
\[
  \{(e_i,f_i,u_i)\}_{i\in I}
\]
where the \(e_i\le e\) are pairwise orthogonal subprojections, the
\(f_i\le f\) are pairwise orthogonal subprojections, and \(u_i\in\cM\) is a
partial isometry with \(u_i^*u_i=e_i\) and \(u_i u_i^*=f_i\).  Order such
families by extension.  Zorn's lemma gives a maximal family.  Since \(\cM\)
is a von Neumann algebra, \texttt{CENTRAL-SUPPORT-001} supplies the
orthogonal projection suprema
\[
  e_0=\bigvee_i e_i,\qquad f_0=\bigvee_i f_i
\]
in \(\cM\), and it supplies the strong sum \(u=\sum_i u_i\) as a partial
isometry in \(\cM\) with \(u^*u=e_0\) and \(uu^*=f_0\).

Let \(r=e-e_0\) and \(s=f-f_0\).  Maximality says that no nonzero
subprojection of \(r\) is equivalent to a subprojection of \(s\).  The
central-support separation field of \texttt{CENTRAL-SUPPORT-001} says that
the residual central supports are orthogonal, so a central projection \(z\)
separates the residual \(e\)-side from the residual \(f\)-side.  The
central-cut Murray--von Neumann field says that if \(a\sim b\) is implemented
by \(v\), then \(za\sim zb\) is implemented by \(zv\).  Applying this to
\(e_0\sim f_0\), and using the central support separation of \(r\) and \(s\),
gives
\[
  ze\precsim zf,\qquad (1-z)f\precsim (1-z)e.
\]
If \(\cM\) is a factor, the only central projections are \(0\) and \(1\).
Thus comparability is global.  Antisymmetry modulo equivalence gives exactly
one of \(e\prec f\), \(e\sim f\), or \(e\succ f\).
\end{proof}

\subsection*{Trace Interfaces}

\begin{definition}[TraceInterfaces, ID: TRACE-001, Takesaki I Definition V.2.1, PDF p.~329]
The positive cone of a concrete von Neumann algebra \(\cM\subseteq\cB(\cH)\) is
\[
  \cM_+=\{x\in\cM:x=x^*\hbox{ and }\langle x\xi,\xi\rangle\ge0
  \hbox{ for all }\xi\in\cH\}.
\]
Equivalently, \(x\in\cM_+\) iff \(x=y^*y\) for some \(y\in\cM\).

A normal positive functional on \(\cM\) is a complex-linear functional
\(\varphi:\cM\to\mathbb C\) such that \(\varphi(x)\ge0\) for all
\(x\in\cM_+\), and for every bounded increasing net
\((x_i)\) in \(\cM_+\) with supremum \(x\in\cM_+\),
\[
  \varphi(x)=\sup_i\varphi(x_i).
\]
It is faithful if \(\varphi(x)=0\) for \(x\in\cM_+\) implies \(x=0\), finite
as a functional when \(\varphi(1)<+\infty\), a state when \(\varphi(1)=1\),
and tracial when
\[
  \varphi(xy)=\varphi(yx),\qquad x,y\in\cM.
\]

An extended positive trace on \(\cM\) is a map
\[
  \tau:\cM_+\to[0,+\infty]
\]
such that
\[
  \tau(x+y)=\tau(x)+\tau(y),\qquad x,y\in\cM_+,
\]
\[
  \tau(\lambda x)=\lambda\tau(x),\qquad \lambda\ge0,\ x\in\cM_+,
\]
with \(0(+\infty)=0\), and
\[
  \tau(x^*x)=\tau(xx^*),\qquad x\in\cM.
\]
It is faithful if \(\tau(x)=0\) for \(x\in\cM_+\) implies \(x=0\), finite if
\(\tau(1)<+\infty\), semifinite if every nonzero \(x\in\cM_+\) majorizes a
nonzero \(y\in\cM_+\) with \(\tau(y)<+\infty\), and normal if
\[
  \tau(\sup_i x_i)=\sup_i\tau(x_i)
\]
for every bounded increasing net in \(\cM_+\).  A finite normal tracial state
\(\varphi\) gives such a trace by restricting \(\varphi\) to \(\cM_+\) and
viewing its nonnegative values in \([0,+\infty]\).
\source{Takesaki I, Definition V.2.1, PDF p.~329.}
\depends{\texttt{VN-002}, \texttt{PROJ-001}.}
\lean{Formalization.OperatorAlgebra.TraceInterfaces}
\status{blueprint_note_complete}
\end{definition}

\subsection*{Type Predicates and Alternatives}

\begin{definition}[FactorTypes, ID: TYPE-DEF-001, Takesaki I Definitions V.1.17, V.1.21, and V.1.33, PDF pp.~316--322]
For a concrete von Neumann algebra \(\cM\), use the projection predicates of
\texttt{PROJ-001} to define:
\begin{itemize}
\item \(\cM\) is type I if every nonzero central projection \(z\in\cM\)
majorizes a nonzero abelian projection \(e\in\cM\), meaning
\[
  0\neq e\le z,\qquad e\cM e\hbox{ is commutative}.
\]
\item \(\cM\) is type III if it has no nonzero finite projection.
\item \(\cM\) is semifinite in the projection/type sense if every nonzero
central projection \(z\in\cM\) majorizes a nonzero finite projection
\(e\in\cM\).  Equivalently, there is no nonzero central projection \(z\) such
that the central summand \(z\cM z\) is type III.
\item \(\cM\) is type II if it is semifinite in this projection/type sense and
has no nonzero abelian projection.
\item \(\cM\) is type \(II_1\) if it is type II and finite, i.e. the identity
projection of \(\cM\) is finite.
\item \(\cM\) is type \(II_\infty\) if it is type II and has no nonzero
central finite projection.
\item A central summand of \(\cM\) is the corner \(z\cM z\) for a central
projection \(z\).  It is nonzero exactly when \(z\neq0\).
\item For the factor-level model predicates used by \texttt{MATRIX-001},
\(\cM\) is type \(I_\alpha\) when it is type I and the identity is
\(\alpha\)-homogeneous: there is an index set \(I\) of cardinal \(\alpha\) and
pairwise orthogonal nonzero abelian projections \((e_i)_{i\in I}\) with
\[
  \bigvee_{i\in I}e_i=1,\qquad e_i\sim_{\cM} e_j\quad(i,j\in I).
\]
For factors this is the cardinal label distinguishing \(I_n\) and
\(I_\infty\).
\end{itemize}
\source{Takesaki I, Definition V.1.17, PDF p.~316; Definition V.1.21, PDF
p.~317; Definition V.1.33, PDF p.~322.}
\depends{\texttt{VN-002}, \texttt{PROJ-001}, \texttt{CENTRAL-SUPPORT-001}.}
\lean{Formalization.OperatorAlgebra.FactorTypes}
\status{blueprint_note_complete}
\end{definition}

\begin{theorem}[SemifiniteIffFaithfulSemifiniteNormalTrace, supporting ID: TRACE-001, Takesaki I Theorem V.2.15, PDF p.~337]
A von Neumann algebra \(\cM\) is semifinite in the projection/type sense of
\texttt{TYPE-DEF-001}, equivalently every nonzero central projection
majorizes a nonzero finite projection, if and only if \(\cM\) admits a
faithful semifinite normal trace.  In particular, a faithful finite normal
trace is a faithful semifinite normal trace.
\source{Takesaki I, Theorem V.2.15, PDF p.~337, and Theorem V.2.4, PDF
p.~331.}
\depends{\texttt{TRACE-001}, \texttt{TYPE-DEF-001},
\texttt{CENTRAL-SUPPORT-001}.}
\lean{Formalization.OperatorAlgebra.semifinite_iff_exists_fsn_trace}
\status{blueprint_note_complete}
\end{theorem}

\begin{proof}
First assume \(\tau\) is a faithful semifinite normal trace.  Let \(z\) be a
nonzero central projection.  Semifiniteness applied to \(z\in\cM_+\) gives a
nonzero \(0<y\le z\) with \(\tau(y)<+\infty\).  Choose a nonzero spectral
projection \(p=1_{[\epsilon,\infty)}(y)\) for some \(\epsilon>0\).  The
functional-calculus input for von Neumann algebras gives \(0\neq p\le z\) in
the corner \(z\cM z\), and \(\tau(p)<+\infty\).  If \(p\) were infinite,
there would be a proper subprojection \(q<p\) with \(q\sim p\).  Traciality
gives \(\tau(q)=\tau(p)\), hence \(\tau(p-q)=0\).  Faithfulness forces
\(p=q\), a contradiction.  Thus every nonzero central projection majorizes a
nonzero finite projection, so \(\cM\) is semifinite in the projection/type
sense.

Conversely assume the projection/type semifiniteness condition.  Takesaki's
central type decomposition splits \(\cM\) into type I, type \(II_1\), and
type \(II_\infty\) central summands, with no type III summand.  On the finite
central part, Theorem V.2.4 supplies finite normal traces whose central
supports cover that summand.  On the properly infinite semifinite type I and
type II parts, the usual construction writes each central piece as a direct
sum or amplification of a finite corner and tensors a faithful finite normal
trace on the finite corner with the standard semifinite trace on the
amplifying \(\cB(\ell^2)\)-part.  Taking the central direct sum of these
traces gives a normal semifinite trace on \(\cM\).  Faithfulness follows
because the supports of the summand traces have central supremum \(1\), using
\texttt{CENTRAL-SUPPORT-001}; if a positive element has trace zero, each
central summand component has trace zero and hence vanishes.
\end{proof}

\begin{theorem}[FactorTypeAlternatives, ID: TYPE-ALT-001, Takesaki I Theorem V.1.19 and Corollary V.1.20, PDF pp.~316--317]
There are unique pairwise orthogonal central projections
\[
  z_I,\ z_{II_1},\ z_{II_\infty},\ z_{III}\in\cZ(\cM)
\]
with sum \(1\) such that the corresponding central summands
\[
  z_I\cM z_I,\quad z_{II_1}\cM z_{II_1},\quad
  z_{II_\infty}\cM z_{II_\infty},\quad z_{III}\cM z_{III}
\]
are respectively type I, type \(II_1\), type \(II_\infty\), and type III,
omitting zero summands.  Consequently, if \(\cM\) is a factor, exactly one of
the following holds:
\[
  \cM\hbox{ is type I},\qquad
  \cM\hbox{ is type }II_1,\qquad
  \cM\hbox{ is type }II_\infty,\qquad
  \cM\hbox{ is type III}.
\]
\source{Takesaki I, Theorem V.1.19 and Corollary V.1.20, PDF pp.~316--317.}
\depends{\texttt{TYPE-DEF-001}, \texttt{CENTRAL-SUPPORT-001}.}
\lean{Formalization.OperatorAlgebra.factor_type_alternatives}
\status{blueprint_note_complete}
\end{theorem}

\begin{proof}
The proof constructs the central summands by central supports.  Let \(z_I\)
be the supremum of the central supports of all abelian projections.  Then
\(z_I\cM z_I\) is type I by construction, while the complement contains no
nonzero abelian projection.  Inside \((1-z_I)\cM(1-z_I)\), let \(z_{II}\) be
the supremum of the central supports of all finite projections.  Then
\(z_{II}\cM z_{II}\) is type II and the complementary summand has no nonzero
finite projection, hence is type III.  Finally split \(z_{II}\) into the
largest central finite projection \(z_{II_1}\) and its complement
\(z_{II_\infty}=z_{II}-z_{II_1}\).  These are central, pairwise orthogonal,
and sum to \(1\).  Each defining maximality statement gives the asserted type
predicate for the corresponding summand.

Uniqueness follows because the properties are detected by central support:
any type I central summand must be below \(z_I\); any finite type II central
summand must be below \(z_{II_1}\); any properly infinite semifinite type II
central summand must be below \(z_{II_\infty}\); and any type III central
summand must be below \(z_{III}\).  If \(\cM\) is a factor, every central
projection is \(0\) or \(1\), so exactly one of the four central summands is
nonzero.  The definitions in \texttt{TYPE-DEF-001} identify that nonzero
summand with the corresponding factor type.
\end{proof}

\subsection*{Represented-Region Vector-State Criteria}

\begin{definition}[RegionVectorState, ID: REGION-VECTOR-001]
Let \(\omega\) be a state on the quasi-local algebra \(\mathcal A_\Gamma\),
let
\[
  (\cH_\omega,\pi_\omega,\Omega_\omega)
\]
be its GNS triple from \texttt{GNS-001}, and let
\[
  \cM_\Lambda(\omega)=\pi_\omega(\mathcal A_\Lambda)''
\]
be the represented region von Neumann algebra from \texttt{VN-001}.  The
represented-region vector state is the functional
\[
  \varphi_{\omega,\Lambda}:\cM_\Lambda(\omega)\to\mathbb C,\qquad
  \varphi_{\omega,\Lambda}(x)=
  \langle x\Omega_\omega,\Omega_\omega\rangle .
\]
Well-formedness is immediate from the inclusion
\(\cM_\Lambda(\omega)\subseteq\cB(\cH_\omega)\): every \(x\in
\cM_\Lambda(\omega)\) acts on \(\Omega_\omega\), and the inner product is a
complex scalar.  The GNS reconstruction formula gives
\[
  \varphi_{\omega,\Lambda}(\pi_\omega(a))=\omega(a),
  \qquad a\in\mathcal A_\Lambda.
\]
\depends{\texttt{GNS-001}, \texttt{VN-001}.}
\lean{Formalization.OperatorAlgebra.RegionVectorState}
\status{blueprint_note_complete}
\end{definition}

\begin{definition}[VectorFunctionalAnalyticFacts, ID: VECTOR-ANALYTIC-001, Takesaki I Section II.1 equation (7), PDF p.~81, vector-functional notation]
The analytic facts used by the represented-region vector-state criteria are
packaged as explicit Lean prerequisites:
\begin{itemize}
\item \(\operatorname{VectorFunctionalNormal}\): for every Hilbert space
\(\cH\) and vector \(\xi\in\cH\), the rank-one vector functional
\[
  \theta_\xi:\cB(\cH)\to\mathbb C,\qquad
  \theta_\xi(T)=\langle T\xi,\xi\rangle
\]
is a normal positive functional in the monotone-positive-net sense of
\texttt{TRACE-001}; that is, if \((T_i)\) is a bounded increasing net in
\(\cB(\cH)_+\) with supremum \(T\), then
\[
  \theta_\xi(T_i)\uparrow\theta_\xi(T).
\]
Its restriction to any concrete von Neumann subalgebra remains normal.
\item \(\operatorname{PositiveSqrtBounded}\): every positive bounded
operator \(x\in\cB(\cH)_+\) has a positive bounded square root \(x^{1/2}\)
with \((x^{1/2})^2=x\).
\item \(\operatorname{PositiveSqrtInVN}\): if \(x\) lies in a concrete von
Neumann algebra \(\cM\), then \(x^{1/2}\in\cM\).  The proof of
\texttt{REGION-VECTOR-FAITH-001} only needs the bounded square root, but the
stronger in-algebra fact is recorded because later projection and trace
proofs often need functional calculus inside \(\cM\).
\end{itemize}
This target is an explicit analytic prerequisite, not a hidden sourced
theorem.  It may later be discharged from mathlib functional calculus and
normal functional APIs, or replaced by exact textbook locators in a future
source-grounding pass.
\source{Explicit project-local analytic prerequisite.  Takesaki I, Section
II.1, equation (7), PDF p.~81, gives the vector-functional notation; exact
normality and square-root source locators are intentionally not claimed here.}
\depends{\texttt{VN-002}, \texttt{TRACE-001}.}
\lean{Formalization.OperatorAlgebra.VectorFunctionalAnalyticFacts}
\status{blueprint_note_complete}
\end{definition}

\begin{proposition}[RegionVectorStateNormal, ID: REGION-VECTOR-NORMAL-001, Takesaki I Section II.1 equation (7), PDF p.~81, vector-functional notation; VECTOR-ANALYTIC-001]
For every state \(\omega\) and region \(\Lambda\), the represented-region
vector state \(\varphi_{\omega,\Lambda}\) is a normal positive state on
\(\cM_\Lambda(\omega)\).
\source{Project-local consequence of \texttt{VECTOR-ANALYTIC-001}.  Takesaki
I, Section II.1, equation (7), PDF p.~81, fixes the vector-functional
notation; normality is supplied explicitly by
\texttt{VectorFunctionalNormal}.}
\depends{\texttt{REGION-VECTOR-001}, \texttt{VECTOR-ANALYTIC-001},
\texttt{TRACE-001}.}
\lean{Formalization.OperatorAlgebra.regionVectorState_normal}
\status{blueprint_note_complete}
\end{proposition}

\begin{proof}
Linearity follows from linearity of the operator action and the inner product.
If \(x\in\cM_\Lambda(\omega)_+\), positivity means
\[
  \langle x\xi,\xi\rangle\ge0,\qquad \xi\in\cH_\omega .
\]
Taking \(\xi=\Omega_\omega\) gives
\[
  \varphi_{\omega,\Lambda}(x)
  =\langle x\Omega_\omega,\Omega_\omega\rangle\ge0.
\]
Since \(\Omega_\omega\) is a unit vector in the GNS construction,
\(\varphi_{\omega,\Lambda}(1)=1\).

For normality, apply the \(\operatorname{VectorFunctionalNormal}\) field of
\texttt{VECTOR-ANALYTIC-001} to the rank-one vector functional
\[
  T\mapsto\langle T\Omega_\omega,\Omega_\omega\rangle
\]
on \(\cB(\cH_\omega)\).  Equivalently, if \((x_i)\) is a bounded increasing
net in \(\cM_\Lambda(\omega)_+\) with supremum \(x\), then monotone
convergence for that normal vector functional gives
\[
  \langle x_i\Omega_\omega,\Omega_\omega\rangle
  \uparrow
  \langle x\Omega_\omega,\Omega_\omega\rangle .
\]
Restricting this normal vector functional from \(\cB(\cH_\omega)\) to the
von Neumann subalgebra \(\cM_\Lambda(\omega)\) gives the stated normality.
\end{proof}

\begin{definition}[RegionVectorTracePredicates, ID: REGION-VECTOR-TRACE-001, Takesaki I Theorem II.2.6 and Theorem II.3.9, PDF pp.~90, 94--95]
For a represented region algebra \(\cM_\Lambda(\omega)\), define
\(\operatorname{RegionVectorTracial}(\omega,\Lambda)\) to mean
\[
  \varphi_{\omega,\Lambda}(xy)=\varphi_{\omega,\Lambda}(yx),
  \qquad x,y\in\cM_\Lambda(\omega).
\]
Define
\(\operatorname{DenseRegionStrongStarTraceHypothesis}(\omega,\Lambda)\) to be
the following Lean-facing predicate with explicit fields:
\begin{itemize}
\item \(\operatorname{denseTracial}\): for all
\(a,b\in\mathcal A_\Lambda\),
\[
  \omega(ab)=\omega(ba),
\]
or equivalently
\[
  \varphi_{\omega,\Lambda}(\pi_\omega(a)\pi_\omega(b))
  =
  \varphi_{\omega,\Lambda}(\pi_\omega(b)\pi_\omega(a)).
\]
\item \(\operatorname{strongStarApprox}\): for every
\(x\in\cM_\Lambda(\omega)\), there is a directed net
\((a_i)_{i\in I_x}\) in \(\mathcal A_\Lambda\) and a constant
\(C_x<+\infty\) such that
\[
  \|\pi_\omega(a_i)\|\le C_x
\]
for all \(i\), and for every vector \(\xi\in\cH_\omega\),
\[
  \pi_\omega(a_i)\xi\to x\xi,
  \qquad
  \pi_\omega(a_i)^*\xi\to x^*\xi .
\]
This is the bounded strongly-star approximating-net field.  In applications
it may be supplied by a Kaplansky-density theorem derived from
\texttt{VN-DCT-001}; in this blueprint it is an explicit hypothesis field.
\item \(\operatorname{uniformBoundControl}\): the preceding bound may be
chosen with \(C_x=\|x\|\) when a norm-preserving Kaplansky-density input is
available; otherwise the predicate stores the finite constant \(C_x\) as part
of the data.
\item \(\operatorname{boundedMultiplicationPassage}\): if a uniformly
bounded net \(T_i\) converges strongly to \(T\), then for every fixed bounded
operator \(S\),
\[
  T_iS\to TS,\qquad ST_i\to ST
\]
strongly whenever the relevant products stay in \(\cB(\cH_\omega)\).  The
same passage is available after taking adjoints when the net is one of the
strong-star approximating nets above.
\item \(\operatorname{vectorFunctionalPassage}\): for every uniformly bounded
strongly convergent net \(T_i\to T\),
\[
  \langle T_i\Omega_\omega,\Omega_\omega\rangle
  \to
  \langle T\Omega_\omega,\Omega_\omega\rangle .
\]
\end{itemize}
This predicate is deliberately stronger than the informal phrase
``tracial on a dense region algebra''.  It records the topology, adjoint
convergence, uniform bounds, and vector-functional limit passage needed for a
Lean proof of traciality on the represented von Neumann algebra.
\source{Takesaki I, Theorem II.2.6, PDF p.~90, for strong/weak/strong-star
continuity on bounded operator sets; Theorem II.3.9, PDF pp.~94--95, for the
double-commutant approximation route recorded in \texttt{VN-DCT-001}.}
\depends{\texttt{REGION-VECTOR-001}, \texttt{TRACE-001},
\texttt{VN-DCT-001}.}
\lean{Formalization.OperatorAlgebra.RegionVectorTracePredicates}
\status{blueprint_note_complete}
\end{definition}

\begin{proposition}[RegionVectorStateTracialFromDenseRegion, supporting ID: REGION-VECTOR-TRACE-001]
Assume \(\operatorname{DenseRegionStrongStarTraceHypothesis}
(\omega,\Lambda)\).  Then
\(\operatorname{RegionVectorTracial}(\omega,\Lambda)\).
\depends{\texttt{REGION-VECTOR-TRACE-001}, \texttt{REGION-VECTOR-NORMAL-001},
\texttt{VN-001}, \texttt{VN-DCT-001}.}
\lean{Formalization.OperatorAlgebra.regionVectorState_tracial_of_dense}
\status{blueprint_note_complete}
\end{proposition}

\begin{proof}
Let \(x,y\in\cM_\Lambda(\omega)\).  Apply the
\(\operatorname{strongStarApprox}\) field to \(x\) and \(y\).  Thus choose
uniformly bounded represented local nets
\(\pi_\omega(a_i)\to x\) and \(\pi_\omega(b_j)\to y\) strongly, with
\(\pi_\omega(a_i)^*\to x^*\) and \(\pi_\omega(b_j)^*\to y^*\) strongly as
well.  The \(\operatorname{boundedMultiplicationPassage}\) field gives, for
fixed \(j\),
\[
  \pi_\omega(a_i)\pi_\omega(b_j)\to x\pi_\omega(b_j)
\]
strongly, and similarly after interchanging the factors.  The
\(\operatorname{vectorFunctionalPassage}\) field then permits passage of
\(\varphi_{\omega,\Lambda}\) through these limits.  For represented local
elements, the \(\operatorname{denseTracial}\) field gives
\[
  \varphi_{\omega,\Lambda}(\pi_\omega(a_i)\pi_\omega(b_j))
  =\omega(a_i b_j)
  =\omega(b_j a_i)
  =\varphi_{\omega,\Lambda}(\pi_\omega(b_j)\pi_\omega(a_i)).
\]
Passing first in \(i\) and then in \(j\) gives
\(\varphi_{\omega,\Lambda}(xy)=\varphi_{\omega,\Lambda}(yx)\).
\end{proof}

\begin{lemma}[RegionVectorStateFaithfulIffSeparating, ID: REGION-VECTOR-FAITH-001, Takesaki I Section II.1 equation (7), PDF p.~81, vector-functional notation; VECTOR-ANALYTIC-001]
Define \(\operatorname{RegionVectorFaithful}(\omega,\Lambda)\) to mean that
for every \(x\in\cM_\Lambda(\omega)_+\),
\[
  \varphi_{\omega,\Lambda}(x)=0\quad\Longrightarrow\quad x=0.
\]
Then \(\operatorname{RegionVectorFaithful}(\omega,\Lambda)\) holds if and
only if \(\Omega_\omega\) is separating for \(\cM_\Lambda(\omega)\), i.e.
\[
  x\Omega_\omega=0,\ x\in\cM_\Lambda(\omega)
  \quad\Longrightarrow\quad x=0.
\]
\source{Takesaki I, Section II.1, equation (7), PDF p.~81, for vector
functionals; positive square roots are supplied by the explicit analytic
prerequisite \texttt{VECTOR-ANALYTIC-001}.}
\depends{\texttt{REGION-VECTOR-001}, \texttt{VECTOR-ANALYTIC-001},
\texttt{TRACE-001}.}
\lean{Formalization.OperatorAlgebra.regionVectorState_faithful_iff_separating}
\status{blueprint_note_complete}
\end{lemma}

\begin{proof}
Assume first that \(\Omega_\omega\) is separating.  If \(0\le x\in
\cM_\Lambda(\omega)\) and \(\varphi_{\omega,\Lambda}(x)=0\), let
\(y=x^{1/2}\) be the positive square root of \(x\) supplied by
\texttt{VECTOR-ANALYTIC-001}.
Then
\[
  0=\varphi_{\omega,\Lambda}(x)
  =\|y\Omega_\omega\|^2,
\]
so \(y\Omega_\omega=0\), and therefore
\[
  x\Omega_\omega=y(y\Omega_\omega)=0.
\]
Since \(x\in\cM_\Lambda(\omega)\), separatingness gives \(x=0\).  Thus the
vector state is faithful.  This direction uses only existence of the bounded
positive square root, not membership of \(y\) in \(\cM_\Lambda(\omega)\).

Conversely assume the vector state is faithful.  If
\(x\Omega_\omega=0\), then
\[
  \varphi_{\omega,\Lambda}(x^*x)=\|x\Omega_\omega\|^2=0.
\]
Faithfulness applied to the positive element \(x^*x\) gives \(x^*x=0\), and
the \(C^*\)-identity gives \(x=0\).  Therefore \(\Omega_\omega\) is
separating.
\end{proof}

\begin{corollary}[FaithfulNormalTracialStateTypeConsequences, ID: TRACE-TYPE-CRITERION-001]
Let \(\cM\subseteq\cB(\cH)\) be a concrete von Neumann algebra.  If
\(\varphi:\cM\to\mathbb C\) is a faithful normal tracial state, then the
positive-cone trace
\[
  \tau_\varphi(x)=\varphi(x),\qquad x\in\cM_+,
\]
viewed as a \([0,+\infty]\)-valued map, is a faithful finite normal trace.
Consequently \(\cM\) is finite, semifinite, and not type III.  If \(\cM\) is
a factor, then the factor alternative reduces to the finite cases: type I or
type \(II_1\).
\depends{\texttt{TRACE-001}, \texttt{TYPE-DEF-001}, \texttt{TYPE-ALT-001}.}
\lean{Formalization.OperatorAlgebra.faithfulNormalTracialState_typeConsequences}
\status{blueprint_note_complete}
\end{corollary}

\begin{proof}
The restriction of \(\varphi\) to \(\cM_+\) is additive and homogeneous
because \(\varphi\) is linear and positive.  Traciality gives
\(\tau_\varphi(x^*x)=\tau_\varphi(xx^*)\), normality is exactly the normality
condition in \texttt{TRACE-001}, faithfulness is inherited from
\(\varphi\), and finiteness follows from \(\tau_\varphi(1)=1\).

To prove finiteness of the algebra, let \(v\in\cM\) satisfy \(v^*v=1\) and
let \(vv^*\le1\) be its range projection.  Traciality gives
\[
  \tau_\varphi(vv^*)=\tau_\varphi(v^*v)=\tau_\varphi(1).
\]
Thus \(\tau_\varphi(1-vv^*)=0\).  Since \(1-vv^*\ge0\) and the trace is
faithful, \(1-vv^*=0\).  Hence every isometry in \(\cM\) is unitary, so the
identity projection is finite.  A faithful finite normal trace is semifinite:
if \(0\neq x\in\cM_+\), take \(y=x\); boundedness gives
\(0\le y\le\|x\|1\), hence
\(\tau_\varphi(y)\le\|x\|\tau_\varphi(1)<+\infty\), and faithfulness gives
\(\tau_\varphi(y)>0\).  Therefore \(\cM\) is semifinite by
\texttt{SemifiniteIffFaithfulSemifiniteNormalTrace}.  Since \(1\) is a
nonzero finite projection, the definition of type III as having no nonzero
finite projection fails.  If \(\cM\) is a factor, \texttt{TYPE-ALT-001}
leaves only type I or type \(II_1\), because type \(II_\infty\) has no
nonzero central finite projection and type III has no nonzero finite
projection.
\end{proof}

\begin{corollary}[RegionVectorStateTypeConsequences, ID: REGION-VECTOR-TYPE-001]
If \(\operatorname{RegionVectorTracial}(\omega,\Lambda)\) and
\(\operatorname{RegionVectorFaithful}(\omega,\Lambda)\), then the represented
region algebra \(\cM_\Lambda(\omega)\) is finite, semifinite, and not type
III.  If \(\cM_\Lambda(\omega)\) is a factor, it is type I or type \(II_1\).
\depends{\texttt{REGION-VECTOR-NORMAL-001}, \texttt{REGION-VECTOR-TRACE-001},
\texttt{REGION-VECTOR-FAITH-001}, \texttt{TRACE-TYPE-CRITERION-001}.}
\lean{Formalization.OperatorAlgebra.regionVectorState_typeConsequences}
\status{blueprint_note_complete}
\end{corollary}

\begin{proof}
By \texttt{REGION-VECTOR-NORMAL-001}, the vector state is normal and positive.
The two hypotheses say exactly that it is tracial and faithful.  Therefore it
is a faithful normal tracial state on \(\cM_\Lambda(\omega)\).  Apply
\texttt{TRACE-TYPE-CRITERION-001}.
\end{proof}

\subsection*{Matrix and Type I Model Factors}

\begin{proposition}[MatrixFactors, ID: MATRIX-001, Takesaki I Proposition V.1.23, Corollary V.1.28, and Definition V.1.33, PDF pp.~318--322]
For finite \(n\), \(M_n(\mathbb C)\) is a finite type \(I_n\) factor.  For an
infinite-dimensional Hilbert space \(\cH\), \(\cB(\cH)\) is a type I factor,
and its type is determined by \(\dim\cH\); in particular \(\cB(\ell^2)\) is
the standard type \(I_\infty\) model factor.
\source{Takesaki I, Proposition V.1.23, PDF p.~318; Corollary V.1.28, PDF
p.~320; Definition V.1.33, PDF p.~322.}
\depends{\texttt{TYPE-DEF-001}.}
\lean{Formalization.OperatorAlgebra.MatrixFactors}
\status{blueprint_note_complete}
\end{proposition}

\begin{proof}
For \(M_n(\mathbb C)\), apply Takesaki Proposition V.1.23 with abelian center
\(\mathbb C\), then use the homogeneous-rank definition of type \(I_n\).  For
\(\cB(\cH)\), apply Corollary V.1.28 and Definition V.1.33.  The case
\(\cH=\ell^2\) gives the countably infinite type I model used later.
\end{proof}

\subsection*{Lean Extraction Notes}

\begin{longtable}{@{}p{0.16\linewidth}p{0.34\linewidth}p{0.40\linewidth}@{}}
\toprule
Target & Intended file & Notes \\
\midrule
\endhead
\texttt{VN-002} &
\texttt{Formalization/OperatorAlgebra/Factor.lean} &
Concrete center and factor predicates for \texttt{VonNeumannAlgebra H}. \\
\texttt{VECTOR-ANALYTIC-001} &
\texttt{Formalization/OperatorAlgebra/VectorFunctionalAnalyticFacts.lean} &
Explicit analytic prerequisites for vector functional normality and positive
square roots of bounded positive operators. \\
\texttt{PROJ-001} &
\texttt{Formalization/OperatorAlgebra/ProjectionComparison.lean} &
Projection-in-\(\cM\), order, central cuts, partial isometries, corners,
finite/infinite/properly infinite predicates, and the comparability proof
route. \\
\texttt{CENTRAL-SUPPORT-001} &
\texttt{Formalization/OperatorAlgebra/CentralSupport.lean} &
Source-backed central-support and projection-lattice package: complete
projection lattice, central supports, central cuts, orthogonal partial-isometry
sums, centrally orthogonal finite/abelian suprema, and largest central finite
projection.  Manifest entry present; target remains pending adversarial review
and formalization. \\
\texttt{TRACE-001} &
\texttt{Formalization/OperatorAlgebra/TraceInterfaces.lean} &
Positive cone, normal positive functionals, extended positive traces, and
faithful finite/semifinite/normal trace predicates; the semifinite trace
recognition theorem is a supporting theorem after \texttt{TYPE-DEF-001}. \\
\texttt{TYPE-DEF-001}, \texttt{TYPE-ALT-001} &
\texttt{Formalization/OperatorAlgebra/FactorTypes.lean} &
Concrete type predicates, semifinite projection/type predicate, central
summands, homogeneous type I labels, and the central-decomposition proof route,
using \texttt{CENTRAL-SUPPORT-001}. \\
\texttt{REGION-VECTOR-001},
\texttt{REGION-VECTOR-NORMAL-001} &
\texttt{Formalization/OperatorAlgebra/RegionVectorState.lean} &
Vector state \(x\mapsto\langle x\Omega,\Omega\rangle\) on represented region
von Neumann algebras and normality as a restricted vector functional. \\
\texttt{REGION-VECTOR-TRACE-001},
\texttt{REGION-VECTOR-FAITH-001} &
\texttt{Formalization/OperatorAlgebra/RegionVectorCriteria.lean} &
Traciality predicate plus explicit dense strong-star trace hypothesis, and a
separate faithfulness predicate with faithfulness iff the GNS vector is
separating. \\
\texttt{TRACE-TYPE-CRITERION-001},
\texttt{REGION-VECTOR-TYPE-001} &
\texttt{Formalization/OperatorAlgebra/TraceTypeCriteria.lean} &
Finite, semifinite, and type III exclusion consequences of a faithful normal
tracial state, specialized to represented region vector states. \\
\texttt{MATRIX-001} &
\texttt{Formalization/OperatorAlgebra/MatrixFactors.lean} &
Finite matrices should use mathlib matrix facts; \(\cB(\ell^2)\) may need a
separable Hilbert-space model wrapper. \\
\bottomrule
\end{longtable}

\begin{thebibliography}{9}
\bibitem{TakesakiI}
M. Takesaki, \emph{Theory of Operator Algebras I}.  Local source:
\path{../sources/takesaki_1.pdf}.
\end{thebibliography}

\end{document}
